% Scope: concepts directly needed for the inspected mutation-effect experiments.
\chapter{THEORETICAL BACKGROUND}\label{chap:background}

\section{Proteins and Their Regulation}
\subsection{Amino Acids, Peptides and Protein Structure}
An amino acid is a molecular building block with a side chain that gives it chemical properties such
as charge or hydrophobicity. Amino acids linked by peptide bonds form a sequence, conventionally
written from the N-terminus to the C-terminus. A \emph{residue} is an amino acid at a particular position
in that sequence. A \emph{peptide} denotes a short chain or protein fragment; the experimental input
in the main study is a 28-residue peptide. A full protein can contain hundreds of residues, fold into
three-dimensional domains, bind partners and undergo chemical modifications. Its linear sequence
and spatial structure therefore describe different aspects of the same molecule.

A \emph{wild-type} (WT) sequence is the specified reference for a comparison; a \emph{mutant} differs
from it. A missense DNA variant changes an encoded amino acid. Here, a single amino acid
substitution is represented directly as $(s,i,a)$, without modelling nucleotide changes, expression or
splicing. Insertions, deletions, truncations and combinations of substitutions have different input
requirements and are outside the evaluated regression task.

\subsection{Stability, Turnover, Degradation and Homeostasis}
Proteostasis describes the coordinated maintenance of protein amount, conformation and location
\citep{balch2008}. \emph{Degradation} is the breakdown of protein molecules. \emph{Turnover}
describes their production and removal over time, and \emph{half-life} is a time required for half of an
initial population to disappear under a stated model and condition. \emph{Abundance} is the amount
present at an observation time; it depends on synthesis as well as removal. None of these quantities is
identical to \emph{thermodynamic folding stability}, which describes the energetic preference of
folded relative to unfolded states. A mutation can alter one of these properties without changing all
of them in the same direction. In particular, the PSI label below is neither a folding free-energy
change nor a half-life expressed in hours.

\subsection{Selective Degradation and Degrons}
In a canonical UPS pathway, E1 activates ubiquitin, E2 carries it and an E3 ligase helps determine
substrate recognition and ubiquitin transfer. Appropriate ubiquitination can lead to substrate
processing by the proteasome. Ubiquitin also has non-degradative roles, so ubiquitination alone is not
proof of degradation. Degrons contribute specificity; N- and C-degrons depend on terminal context,
while internal degrons occur within a sequence \citep{varshavsky2019,koren2018}. The UPS is one
part of cellular degradation, alongside pathways such as lysosomal degradation.

Degron activity depends on more than the presence of a short word in an amino acid string. A motif
may be masked by folding or a binding partner, require a post-translational modification, or be inactive
when its recognising machinery is absent. Molecular recognition, access to an appropriate
ubiquitination site and proteasomal processing are distinct requirements. The schematic in
Figure~\ref{fig:ups-concept} illustrates a possible route rather than a universal mechanism.

\begin{figure}[htbp]
\centering
\begin{tikzpicture}[node distance=6mm, every node/.style={font=\small,align=center},
 box/.style={draw,rounded corners,text width=10.4cm,inner sep=5pt},
 arr/.style={-{Latex},thick}]
\node[box] (signal) {Protein with an accessible degradation signal};
\node[box,below=of signal] (recognition) {Recognition in the relevant cellular context\\
E3-mediated substrate selection; E1/E2/E3 ubiquitination machinery};
\node[box,below=of recognition] (removal) {Suitable ubiquitin signal and proteasomal engagement\\
Unfolding and degradation of the substrate};
\draw[arr] (signal) -- (recognition);
\draw[arr] (recognition) -- (removal);
\end{tikzpicture}
\caption{Original conceptual schematic of a canonical UPS route, informed by
\citep{varshavsky2019,koren2018}. It is not an observed CT553H mechanism; individual signals can
require additional conditions, and not all ubiquitination leads to this route.}
\label{fig:ups-concept}
\end{figure}

\subsection{How a Substitution Can Change a Degradation Signal}
A replacement can remove a favourable recognition residue, create a new recognition sequence,
alter local chemistry or change how a region is exposed. These are alternative mechanisms, not
mechanistic labels supplied by the current model. Degron annotations and experimentally investigated
variant effects support the relevance of both sequence and structural context
\citep{degronopedia,degronmd,voutsinos2026}. Conversely, a mutation can alter catalytic activity,
localisation or binding without acting through a degron. A measured reporter effect therefore warrants
a statement about that reporter endpoint, followed by a separately labelled biological interpretation.

\section{Experimental Measurements and Mutation Labels}
\subsection{GPS Profiling and the Meaning of PSI}
Global Protein Stability profiling uses a fluorescent fusion and a reference fluorescent signal to
compare reporter abundance in cells \citep{yen2008}. In the Zhang implementation, GFP is fused to a
peptide and fluorescence relative to DsRed determines six sorting bins. Sequencing associates
constructs with their distribution across these bins \citep{zhang2023}.

% Source of exact reconstruction: research/full-data-audit/build_psi.py, build().
For a construct $v$, let $c_{vb}$ be its read count in bin $b$, and let $C_b$ be the total count used to
normalise that bin in the reconstruction. The inspected processing code uses
\begin{equation}
 n_{vb}=c_{vb}\frac{C_1}{C_b},\qquad
 \mathrm{PSI}_v=\frac{\sum_{b=1}^{6}b\,n_{vb}}{\sum_{b=1}^{6}n_{vb}}.
 \label{eq:psi-bins}
\end{equation}
PSI is defined only with sufficient nonzero read support. Its scale is 1--6: lower values correspond
to less stable GFP--peptide reporters and higher values to more stable reporters under this assay.
The normalisation population and handling of headers matter when reconstructing it from counts.
It is a bin-based index, not a direct degradation-rate measurement. GFP maturation, fusion design,
expression and cellular condition are part of the observation process. The model does not remove
those factors merely by predicting the index accurately.

\subsection{Scanning, Saturation and the Unit of Observation}
Scanning assays test selected replacements along many peptides. Saturation assays seek many or
all alternative amino acids at each position, trading breadth across peptides for depth within each
peptide. A complete single-substitution design has $19L$ non-self replacements for length $L$;
actual observation is limited by library construction, read coverage and quality control. Systematic
sequence--function mapping therefore records both observed effects and missing variants
\citep{fowler2010}.

In this report a \emph{mutation cell} means one entry $(\text{peptide},i,a)$ in a mutation matrix; it is
not one biological cell. Many mutation cells share a peptide, and several peptides can share a source
protein or a homologous protein group. They share reference measurements and sequence context.
Consequently, thousands of mutation cells do not provide thousands of independent protein-level
replications. Experimental replicates and random-seed training repeats are also different: the former
probe measurement variation, while the latter probe variation in the computational procedure.

\subsection{Constructing and Interpreting the Effect}
For each technically eligible mutation, the label is the mutant PSI minus a compatible WT reference.
% Source: residue-task-derisk-2026-10-15/residue_data.py, build_saturation().
In the inspected saturation processing, each of the two measurement halves has its own WT baseline,
obtained as the median of at least ten measured entries encoding the WT amino acid. The final label
averages the two within-half differences. Scanning uses the directly measured WT construct for its
condition. An unrelated BAG6-WT condition is not substituted for these references.
\begin{equation}
 \Delta\mathrm{PSI}_{p,i,a}
 =\frac{1}{2}\sum_{r\in\{A,B\}}
       \bigl(\mathrm{PSI}_{p,i,a,r}-\mathrm{PSI}_{p,\mathrm{WT},r}\bigr).
 \label{eq:delta-psi-halves}
\end{equation}
A positive effect means reporter stabilisation; a negative effect means destabilisation. Interpretation
as loss or creation of a degron remains conditional on the experimental system. The WT reference is
part of label construction, not an experimentally available feature required at prediction time.
Figure~\ref{fig:mutation-observation} separates this measurement route from the prediction route.

\begin{figure}[htbp]
\centering
\begin{tikzpicture}[every node/.style={font=\small,align=center},
 box/.style={draw,rounded corners,text width=5.1cm,inner sep=5pt}, arr/.style={-{Latex},thick}]
\node[box] (input) at (0,0) {WT sequence, position, replacement};
\node[box] (lab) at (6.5,0) {WT and mutant reporter measurements};
\node[box] (prediction) at (0,-1.8) {Sequence model\\Predicted $\widehat{\Delta\mathrm{PSI}}$};
\node[box] (label) at (6.5,-1.8) {Compatible reference subtraction\\Observed $\Delta\mathrm{PSI}$};
\draw[arr] (input) -- (prediction);
\draw[arr] (lab) -- (label);
\draw[arr] (prediction) -- node[above,text width=1.2cm] {Error} (label);
\end{tikzpicture}
\caption{Original conceptual diagram of mutation-effect prediction and label construction. No
hypothetical PSI values are presented as observations. Measured WT PSI belongs to the label route.}
\label{fig:mutation-observation}
\end{figure}

\subsection{Quality Control, Noise and PSI-Based Selection}
Read counts affect precision; low support, invalid coordinates or residues, missing WT references and
unreliable replicate support can prevent a label from being usable. Disagreement between replicates
can reflect technical noise or biological condition differences. Replicate averaging reduces some
variation but does not establish that errors are independent or remove systematic biases. Abundance
assays illustrate the need to retain replicate-level uncertainty and assay applicability limits
\citep{matreyek2018}. The exact filters and their sample-count effects belong to the data methods,
rather than to an assumed universal definition of a degron.

WT-PSI thresholds select different populations. Because both mutant and reference PSI lie between
1 and 6, an effect is bounded by $1-\mathrm{PSI}_{\mathrm{WT}}$ and
$6-\mathrm{PSI}_{\mathrm{WT}}$. A high-PSI WT has less room for large positive changes; the event
prevalence and regression target distribution can therefore differ between strata. Expanding the
training range changes which examples the model learns from, not just the number of rows.
In v3, WT PSI defines training arms and reporting strata; counts support QC, and disagreement is used
only in a designated training-weight experiment. WT PSI, counts, groups and fold identifiers are not
sequence-model inputs. Test predictions are not clipped using the measured WT PSI.

\section{Representations and Learning Methods}
\subsection{Supervised Regression and Simple Controls}
Supervised regression learns a mapping from inputs to a continuous target using paired observations.
A zero-effect predictor provides a necessary reference when many substitutions have small effects.
One-hot encoding represents an amino acid by a vector with one nonzero entry among 20 categories;
it does not assume a numerical order among amino acids. Peptide composition records counts or
fractions and discards order. Ordered local windows preserve which residue occurs at each relative
position, allowing a comparison between context and composition.

Ridge regression adds an $L_2$ penalty to squared-error fitting \citep{hoerl1970}:
\begin{equation}
 \widehat{\beta}=\arg\min_{\beta}
 \sum_j(y_j-x_j^\mathsf{T}\beta)^2+\lambda\|\beta\|_2^2.
\end{equation}
The penalty stabilises estimation with correlated features. Ridge is inexpensive and interpretable at
the feature level, but its fitted response is linear in the chosen representation. Interactions require
explicit features or a nonlinear model. Feature scaling and the penalty must be chosen without using
outer-test observations.

\subsection{Protein Language Models and Embeddings}
A pLM learns sequence regularities from protein corpora. ESM-2 uses a transformer to learn contextual
residue representations \citep{lin2023esm,vaswani2017}. An \emph{embedding} is a vector representation,
not a directly measured biological property. A residue vector can depend on the rest of the input
sequence. The current frozen-feature design combines the change in the substituted-site vector and
the change in mean sequence vectors, optionally with ordered context. A trained head maps these
features to $\Delta$PSI. Dimensionality reduction limits feature size; its parameters must be fitted on
training data within each selection split.

Pretraining can provide useful regularities when endpoint labels are limited, but usefulness must be
measured against simple controls. Larger embeddings or checkpoints need not improve this endpoint.
The labelled-train/test separation also does not certify that sequences or homologs were absent from
the original pLM corpus. Sequence pretraining and supervised endpoint-label leakage require separate
audits; neither is dismissed by a claim that the encoder is frozen.

\subsection{Pointwise, CNN and Multi-Scale CNN Heads}
A one-dimensional convolution applies a learned filter along sequence positions. A kernel of width
$k$ mixes information across $k$ neighbouring positions; applying the same filter at each position
shares parameters. The pointwise $k=1$ head is a control without additional neighbouring-position
mixing at that layer. On contextual pLM embeddings, it can still use context already encoded by the
pLM. A CNN with a wider kernel tests additional local mixing; parallel kernels of different widths form
a multi-scale CNN (mCNN). These are inductive biases, not proof of a biological motif mechanism.

In v3 the tested heads include $k=1$, $k=5$ and parallel $k=3,5,7$ branches, with shared centre and
sequence-pair/composition inputs where specified. This tests the added role of mixing while keeping
important inputs comparable. It does not make all heads identical in parameter count. Additional
capacity can fit noise; early stopping, matched folds and paired uncertainty help assess whether the
extra complexity has reproducible predictive value. PAP is a related CNN-based peptide-abundance
predictor; a mutant--WT PAP difference is a reference score, not automatically a calibrated PSI effect
\citep{voutsinos2026}.

\subsection{Frozen Encoders, Fine-Tuning and LoRA}
A frozen encoder retains its pretrained weights while the regression head is learned. Full fine-tuning
updates the encoder as well. LoRA instead learns a low-rank update, schematically
$W'=W+BA$, while holding the base matrix $W$ fixed \citep{hu2022lora}. This reduces the number of
trainable adaptation parameters; it does not guarantee lower total runtime or greater accuracy for this
dataset. The implemented comparison adds LoRA to attention query/value projections and learns a
head, using a frozen-encoder control with the same data assignments and base revision. Observing
changed adapter parameters establishes that adaptation occurred; completed held-out predictions are
still required to establish its performance. LoRA evidence cannot stand in for an uncompleted
full-fine-tuning experiment.

\subsection{Transfer Learning and Endpoint Compatibility}
Transfer learning uses information learned on a source task or population for a target task. Scanning
mutation labels can train a mutation-effect representation with narrower amino-acid coverage;
peptide-abundance pre-training supplies a different endpoint and typically a different head. The
source must exclude prohibited overlap with target-test groups before training. Source similarity,
assay differences and target-label coverage can make transfer helpful or harmful. Negative transfer
is an empirical result to retain, not evidence that an implementation was necessarily defective.
Binary motif annotations are not converted into continuous $\Delta$PSI labels.

\section{Evaluation and Reliability}
\subsection{Data Splitting, Leakage and Generalisation}
Leakage occurs when information unavailable in the intended prediction setting influences training or
selection, yielding misleading evaluation \citep{bernett2024}. Randomly splitting mutation cells can
put nearly identical mutants and shared WT sequences on both sides. Protein-group-aware splitting
keeps related source proteins together; MMseqs2 supports the sequence comparisons
\citep{steinegger2017}. The historical cov80/cov50 definitions and v3 unions are described separately
in the methods. A heuristic similarity threshold reduces specified overlap but does not prove complete
biological independence, especially for short shared motifs or remote homology.

Nested evaluation chooses hyperparameters, scalers, PCA and stopping epochs inside outer-training
data, then evaluates on held-out outer groups. An out-of-fold (OOF) prediction is made by a fitted
model that excluded the corresponding outer group. Overfitting can still occur through repeated
human selection on the same development OOF results. A locked recipe evaluated on an untouched
source provides different evidence from selecting a winner on development data; a source previously
explored in historical work must be described accordingly.

\subsection{Losses and Regression Metrics}
The inspected regressors use squared-error objectives or inner-validation MSE for selection. Selection
loss and reporting metric need not coincide. For the main full-PSI endpoint, peptide-macro MAE is
\begin{equation}
 \mathrm{MAE}_{\mathrm{macro}}=
 \frac{1}{P}\sum_{p=1}^{P}\frac{1}{n_p}\sum_{j=1}^{n_p}
       |y_{pj}-\widehat{y}_{pj}|.
 \label{eq:macro-mae}
\end{equation}
This gives each eligible peptide equal weight, rather than giving a densely observed peptide more
influence solely because it has more mutation cells. RMSE increases the influence of large errors;
Spearman correlation measures rank association rather than error magnitude. Their aggregation and
eligibility rules must be stated before comparing values.

\subsection{Event Ranking and Probability Calibration}
The v3 positive and negative events are $\Delta\mathrm{PSI}\geq1$ and
$\Delta\mathrm{PSI}\leq-1$, respectively. AUROC measures ranking of event versus non-event scores;
average precision and top-fraction precision depend additionally on event prevalence. Peptides
without sufficient examples of both classes have undefined event-ranking metrics and remain
eligible for regression. The old restricted benchmark uses its own eligibility and macro-AUROC rules.

A regression score is not a probability. Calibration learns event probabilities using permitted
validation predictions. The Brier score measures squared probability error, whereas expected
calibration error (ECE) summarises discrepancies between predicted and observed frequencies within
chosen bins. Bin choice and population composition affect ECE. Improved ranking, Brier score and ECE
need not occur together, and development calibration does not guarantee calibration after a source
shift. A calibrator must not see the outer-test outcomes it is later evaluated on.

\subsection{Uncertainty, Seeds and Fair Comparisons}
A paired comparison evaluates both methods on the same eligible observations. Bootstrap intervals
approximate variation under resampling \citep{efron1979}; here the resampling unit is the protein
group, preserving the dependence of its mutation cells. The interval addresses uncertainty under this
sampling scheme, not unmeasured assay biases or all future cell types. Seeds measure computational
variation and are reported alongside ensemble predictions. The metric of an averaged prediction is
not generally equal to the average of seed-specific metrics.

For a difference defined as MAE(method)--MAE(control), a negative value favours the method. A
confidence interval crossing zero leaves the direction unresolved at the stated level; it does not
prove equivalence. Multiplicity adjustment, when specified, applies to the declared comparison family,
not every exploratory contrast. Practical relevance also requires effect size, stratum behaviour and
cost. Comparisons across different endpoints, QC versions, folds or runtime campaigns are contextual
unless those differences have been reconciled explicitly.

\section{Research and Software Toolchains}
Python, NumPy, pandas, scikit-learn and PyTorch implement the research pipeline. Pinned revisions,
configuration, manifests and prediction artifacts support reconstruction. React/TypeScript, FastAPI,
local persistence and an inference worker implement the application design. Software test success
checks a stated behaviour; it does not supply biological validation, independent model generalisation
or proof of wet-lab utility.
